\section{Erzählverhalten}

\startplacetable[location={here,none},title={}]
\setupTABLE[c][2][loffset=0.5cm]
\bTABLE[frame=off]
\bTR 
	\bTD [background=color,backgroundcolor=gray] Erzählvorm: \eTD 
	\bTD [nr=2]
\startitemize[packed]
\item 1	
\item 2	
\item 3	
\stopitemize
	\eTD
\eTR
\bTR
	\bTD Ich-Erzählung, Perpektive der Tochter \eTD 
\eTR
\bTR 
	\bTD [background=color,backgroundcolor=gray] Erzählperspektive: \eTD 
	\bTD [nr=2]
\startitemize[packed]
\item 1	
\item 2	
\item 3	
\stopitemize
	\eTD
\eTR
\bTR
	\bTD Außensicht \eTD 
\eTR
\bTR 
	\bTD [background=color,backgroundcolor=gray] Erzählverhalten: \eTD 
	\bTD [nr=2]
\startitemize[packed]
\item 1	
\item 2	
\item 3	
\stopitemize
	\eTD
\eTR
\bTR
	\bTD Personal \eTD 
\eTR
\bTR 
	\bTD [background=color,backgroundcolor=gray] Erzählerstandort: \eTD 
	\bTD [nr=2]
\startitemize[packed]
\item 1	
\item 2	
\item 3	
\stopitemize
	\eTD
\eTR
\bTR
	\bTD Teil des Geschehen \eTD 
\eTR
\bTR 
	\bTD [background=color,backgroundcolor=gray] Erzählhaltung: \eTD 
	\bTD [nr=2]
\startitemize[packed]
\item 1	
\item 2	
\item 3	
\stopitemize
	\eTD

\eTR
\bTR
	\bTD wertend \eTD 
\eTR
\eTABLE
\stopplacetable

\subsection{Funktion der Muschel}

\startitemize[packed]
\item {\bf Schutz und Verteidigung} vor natürlichen Feinden aufgrund:
\startitemize[packed]
\item ihre Beschaffenheit: harte Schale, weiches Inneres
\item Funktion: geschlossene Muschelschalen bilden eine Art Gefäß, Symbol für sensible, leicht verletzbare, empfindliche Persönlichkeit
\stopitemize

\item {\bf Schönheit und Anmut} aufgrund ihrer
\startitemize[packed]
\item runde Form
\item konzentrisch angeordneten, feinen Schüppchenreihen
\stopitemize

\item {\bf Fruchtbarkeit}
\startitemize[packed]
\item Heranreifen einer wertvollen zw. Kostbaren Perle im Inneren
\item Entstehung neuen Lebens, Symbol für das weibliche Geschlechtsorgan und lust.
\stopitemize

\item {\bf Aphrodisiakum} zur Belebung der Libido bzw. des sexuellen Lustempfindens
\item Spritiuell: {\bf Zyklus} von Leben und Tod
\startitemize[packed]
\item Reife und Vervollkommnung der Perle im Verborgenen, ehe sie das Licht der Welt erblickt
\item Perle überdauert mehrere Generationen, während die Muschel selbst längst tot ist. 
\startitemize[packed]
\item Symbol für das Wunder im Meer der Zeit
\stopitemize
\stopitemize

\item Im Christentum: {\bf Auferstehung}
\startitemize[packed]
\item Beerdigung und anschließende Auferstehung Jesus
\item gewaltsame Öffnung des fast verschlossenen Grabes
\item Mittelalter: Pilger-Jakobsmuschel als Schmuck der Gewänder von Wallfahrern
\startitemize[packed]
\item Symbol für das Grab des Herrn als kostbare Perle
\stopitemize
\stopitemize

\stopitemize

\startframedtask
{\bf Arbeiten} Sie mithilfer der folgenden Textstellen die verschiednene Assoziationen der Figuren beim Thema "Muscheln" {\bf heaus}.

\blank

\bTABLE
\startCSV
Vater und Mutter:, "S. 8, 10"
Tochter:, "S. 10, 18-21, 35, 48f., 128."
\stopCSV
\eTABLE
\stopframedtask

\startitemize[packed]
\item Vater und Sohn mögen als einizege Muscheln
\item Die Mutter mag keine Muscheln, und möchte auch diese nicht zubereiten, tut es aber trozdem.
\stopitemize
